\documentclass[12pt,a4paper]{article}
\usepackage[utf8]{inputenc}
\usepackage[T1]{fontenc}
\usepackage[french]{babel}
\usepackage{xcolor}
\usepackage{geometry}
\geometry{a4paper, margin=2cm}
\usepackage{enumitem}
\usepackage{hyperref}
\usepackage{sectsty}

\definecolor{primary}{RGB}{33, 100, 160}
\definecolor{secondary}{RGB}{200, 50, 50}
\definecolor{lightbg}{RGB}{245, 245, 250}

\sectionfont{\color{primary}}
\subsectionfont{\color{secondary}}

\title{\textbf{\Huge Le Rôle du SLA}\\ \large \vspace{0.5cm} \textit{Projet IncidentFlow}}
\author{}
\date{}

\begin{document}
\maketitle
\vspace{-1.5cm}
\hrule
\vspace{1cm}

\section*{Qu'est-ce qu'un SLA (Service Level Agreement) ?}
Un \textbf{SLA} (Accord de Niveau de Service) est un engagement entre un fournisseur de service et son client. Il définit précisément le \textbf{niveau de qualité attendu} (comme le temps de réponse ou de résolution) et ce qui se passe si ce niveau n'est pas atteint.

\vspace{0.5cm}
\noindent\fcolorbox{primary}{lightbg}{
\begin{minipage}{0.97\textwidth}
\textbf{Astuce - L'Analogie de la livraison de Pizza :}\\
Si vous commandez une pizza avec la promesse d'une livraison en moins de 30 minutes, sinon elle est gratuite.
\begin{itemize}[label=--]
    \item \textbf{Le service :} La livraison de la pizza.
    \item \textbf{La métrique :} Le temps (moins de 30 minutes).
    \item \textbf{La pénalité :} La pizza est offerte en cas de retard.
\end{itemize}
\end{minipage}}

\vspace{1cm}

\section*{Le Rôle du SLA dans IncidentFlow}
Dans l'application \textbf{IncidentFlow}, le SLA n'est pas qu'un simple indicateur statistique. C'est le \textbf{moteur de l'urgence} qui garantit la réactivité du support technique. Il est défini par la variable \texttt{sla\_due\_at} et surveillé activement par le composant \texttt{EscalationService}.

\subsection*{1. Fixer une date limite stricte}
Chaque incident possède un champ \texttt{sla\_due\_at} qui représente l'échéance accordée à l'équipe technique pour prendre en charge ou résoudre le problème avant d'être considéré en retard.

\subsection*{2. Surveillance continue}
Le système intègre un planificateur automatique qui vérifie en permanence l'état des délais. Il réagit dans deux situations critiques :
\begin{itemize}[label=\textcolor{primary}{\textbullet}]
    \item \textbf{La Menace d'expiration :} Il reste moins de 15 minutes avant la date limite SLA et l'incident est toujours inactif (au statut "Nouveau").
    \item \textbf{Le Dépassement avéré :} La date limite SLA est officiellement dépassée.
\end{itemize}

\subsection*{3. Le déclenchement de l'Escalade Automatique}
Lorsqu'un SLA est menacé ou violé, IncidentFlow déclenche instantanément une procédure de crise (l'escalade) comprenant 5 actions correctives :
\begin{enumerate}[label=\textbf{\arabic*.}]
    \item \textbf{Montée en priorité :} La priorité de l'incident est forcée à \textbf{"Critical"}.
    \item \textbf{Réassignation immédiate :} L'incident est retiré à son responsable actuel et transféré automatiquement à un supérieur (\textit{Responsable} ou \textit{Administrateur}).
    \item \textbf{Marquage du ticket :} L'incident est officiellement marqué avec le statut \texttt{Escalated}.
    \item \textbf{Traçabilité totale :} Le système injecte un commentaire d'alerte dans le fil de discussion du ticket et trace cette action dans l'historique d'audit (IncidentHistory).
    \item \textbf{Notification :} Une alerte est simulée/envoyée au nouveau manager pour l'informer de la situation d'urgence.
\end{enumerate}

\vspace{1.5cm}
\begin{center}
    \textit{Document généré pour la documentation et la présentation du projet IncidentFlow.}
\end{center}

\end{document}
